% ============================================================================
% Chapter 16 --- Conformance and replaceability tests
% ----------------------------------------------------------------------------
% Purpose : operationalise the interface contracts of Ch. 6--14 and the four
%           replaceability conditions of Ch. 5 through a structured test
%           suite. Closes Part II.
% Page target: ~3 pages.
% Dependencies: Ch. 5 (replaceability conditions), Ch. 6--14 (each layer's
%               contract), Ch. 18--27 (each wave closes via this suite).
% ----------------------------------------------------------------------------
% Decisions registered in _coherence-EN.md:
%   D16.1 four test categories: interface fidelity, data portability,
%         exit rehearsal, decision auditability
%   D16.2 reference test suite across the nine layers
%   D16.3 conformance day procedure
%   D16.4 reporting template with exit-confidence index per layer
%   D16.5 each wave in Part III closes by passing the relevant subset
% ============================================================================

\chapter{Conformance and replaceability tests}
\label{ch:conformance}

The contracts of Chapters~6--14 are stated as obligations on
component behaviour; the conditions of replaceability of
\trchap{ch:architectural-principles}{} (\S\ref{sec:ap-replaceability}) are stated as
properties. Without a mechanism to verify that the obligations are
met and the properties hold, both reduce to declarations. This
chapter operationalises them through a structured test suite that
the institution exercises at agreed cadences and at every wave
closure of Part~III.

%-----------------------------------------------------------------------------
\section{Why a conformance suite}
\label{sec:cf-why}

A contract that cannot be checked is a slogan. The conformance suite
is the artefact that turns each layer's contract from a procurement
clause into a verifiable property of the running infrastructure.
Three uses follow:

\begin{itemize}[leftmargin=1.5em,nosep]
  \item \textbf{Acceptance.} A component is acceptable into
    production when it passes the tests applicable to its layer.
  \item \textbf{Drift detection.} A component that previously
    passed and now fails reveals operational or configuration
    drift before the drift becomes incident.
  \item \textbf{Wave closure.} A migration wave (Part~III) closes
    by passing the relevant subset of the suite, not by the
    subjective judgement that ``the migration is done''.
\end{itemize}

The suite does not replace security testing, performance testing,
or functional testing of the institution's applications. It tests
that the architecture's contracts and replaceability properties
hold, which is a different and structurally narrower question.

%-----------------------------------------------------------------------------
\section{Test categories}
\label{sec:cf-categories}

Tests are organised in four categories.

\begin{description}[leftmargin=1.5em,style=nextline,nosep]
  \item[Interface fidelity.] The component speaks the standard
    interface correctly. The test exercises the standard's defined
    behaviour and verifies the response.
  \item[Data portability.] The institution's data and configuration
    can be exported in a documented open format and re-imported into
    an alternative component without loss.
  \item[Exit rehearsal.] The replacement of the component by an
    alternative can be performed within a documented time window,
    with the institution's data preserved and its operations
    sustained. The test is performed on a representative subset; it
    is rehearsed, not merely simulated.
  \item[Decision auditability.] The events and decisions emitted
    by the component are sufficient to reconstruct any consequential
    action. The test takes a recorded decision and verifies that it
    can be replayed against the recorded inputs and policy version
    to produce the same result.
\end{description}

\noindent The four categories together correspond to the four
conditions of replaceability of \S\ref{sec:ap-replaceability}. A layer that scores well on
fidelity but poorly on exit rehearsal is, despite appearances,
captive.

%-----------------------------------------------------------------------------
\section{Reference test suite per layer}
\label{sec:cf-suite}

The reference suite below organises the conformance tests
in four categories across the nine layers. The suite is illustrative and is intended to
seed the institution's own; tests should be added or removed as the
institutional context warrants.

\begin{table}[H]
\centering\small
\caption{Reference test suite (illustrative). Categories: F~=~Interface
fidelity, P~=~Data portability, X~=~Exit rehearsal, A~=~Decision
auditability.}
\label{tab:cf-suite}
\begin{tabularx}{\textwidth}{%
  >{\hsize=0.6\hsize\linewidth=\hsize\bfseries\raggedright\arraybackslash}X%
  >{\hsize=0.25\hsize\linewidth=\hsize\centering\arraybackslash}X%
  >{\hsize=2.15\hsize\linewidth=\hsize\raggedright\arraybackslash}X%
}
\toprule
Layer & Cat. & Test \\
\midrule
\multirow{4}{*}{Identity (Ch. 6)}
  & F & OIDC discovery endpoint returns metadata conformant to OpenID Connect Discovery 1.0. \\
  & F & Authentication assertion carries the institutional attribute schema, with no captive extensions in the assertion path. \\
  & P & SCIM 2.0 user export and re-import on an alternative IdP preserves attributes and group memberships. \\
  & A & Authentication and session-lifecycle events allow forensic reconstruction of any session. \\
\addlinespace
\multirow{4}{*}{Policy (Ch. 7)}
  & F & PDP query API returns decisions in the canonical \texttt{\{verdict, obligations, policy\_version\}} shape. \\
  & A & Recorded decision replayed against recorded inputs and policy version produces an identical verdict. \\
  & P & Policy bundle exported, re-imported into an alternative engine, produces identical decisions on a representative input set. \\
  & F & Shadow-mode decisions are logged with the same fidelity as enforcing-mode decisions. \\
\addlinespace
\multirow{4}{*}{Network (Ch. 8)}
  & F & NETCONF/YANG configuration apply succeeds without recourse to the supplier's management console (D8.7). \\
  & F & RADIUS attachment yields the expected zone assignment given a representative subject. \\
  & F & IPFIX flow records carry the canonical \texttt{correlation\_id}. \\
  & X & Configuration exported and partially re-imported on an alternative-vendor switch restores expected connectivity on a test segment. \\
\addlinespace
\multirow{3}{*}{Compute (Ch. 9)}
  & F & Kubernetes API passes upstream conformance tests on the production cluster. \\
  & P & Container image migration from the production registry to an alternative OCI-compliant registry preserves the image digest. \\
  & A & SLSA attestations verify against expected provenance for a representative artefact. \\
\addlinespace
\multirow{4}{*}{Storage (Ch. 10)}
  & F & S3 API operations on the production object store interoperate with an independent S3-compatible implementation on a representative test set. \\
  & P & KMIP key migration to an alternative key manager preserves the key material's usability. \\
  & A & Snapshot inventory matches the published retention policy without orphans. \\
  & P & Classification labels preserved across an export-import cycle. \\
\addlinespace
\multirow{3}{*}{Endpoint (Ch. 11)}
  & F & Posture data from each managed OS arrives at the policy layer in the canonical schema. \\
  & F & Decommissioning of a representative endpoint executes the documented procedure end-to-end. \\
  & P & MDM configuration profiles export to a documented format that an alternative MDM can consume on a representative subset. \\
\addlinespace
\multirow{4}{*}{Observability (Ch. 12)}
  & F & Events from every layer arrive at the analytical store conformant to the canonical TAC schema. \\
  & F & Retention horizon verified by deliberate retrieval at the documented limit. \\
  & A & Cross-layer correlation via \texttt{correlation\_id} succeeds for a representative TAC activation. \\
  & X & Event re-routing from the primary analytical store to an alternative produces the same query results on a sampling window. \\
\addlinespace
\multirow{3}{*}{Crypto (Ch. 13)}
  & F & ACME issuance conforms to RFC 8555 against the institutional CA. \\
  & A & Key-escrow recovery procedure executes successfully on a representative escrowed key. \\
  & X & HSM key migration to an alternative HSM preserves the key's cryptographic identity on a representative subset; for non-exportable trust-anchor keys, the key-rotation / re-anchoring rehearsal declared in the Ch.~13 contract executes under dual control. \\
\addlinespace
\multirow{3}{*}{Orchestration (Ch. 14)}
  & F & Identical TAC activation inputs, replayed, produce identical end state (idempotency). \\
  & A & TAC deactivation reverses every consequential change recorded under the activation. \\
  & X & Workflow definitions exported and re-imported on an alternative engine reproduce the same outcome on a representative case. \\
\bottomrule
\end{tabularx}
\end{table}

%-----------------------------------------------------------------------------
\section{How to run a conformance day}
\label{sec:cf-day}

A \emph{conformance day} is a scheduled, time-boxed exercise of the
suite, conducted by a cross-functional team. The reference structure
is:

\begin{itemize}[leftmargin=1.5em,nosep]
  \item \textbf{Scope.} The subset of tests applicable to the
    current scope (full institutional review, single layer, single
    wave closure).
  \item \textbf{Participants.} One engineer per affected layer,
    one observability operator (to retrieve replay material), one
    note-taker, and the operational owner of the scope.
  \item \textbf{Duration.} Indicatively, half a day for a single
    layer; one day for a full institutional review.
  \item \textbf{Output.} A structured report (\S\ref{sec:cf-reporting})
    and a remediation timeline for any test that did not pass.
\end{itemize}

\paragraph{Cadence.} For the layers the institution operates itself,
the full applicable suite is exercised at a regular cadence (a yearly
minimum is a reasonable default for a research-intensive operator; a
down-scaled adopter sets the cadence in proportion to the layers it
self-operates). Layer-specific subsets are exercised at the closure of each
relevant wave (Part~III), and ad hoc when a component is replaced or
when its supplier announces a material change.

%-----------------------------------------------------------------------------
\section{Reporting template}
\label{sec:cf-reporting}

The conformance report records, for each test exercised:

\begin{itemize}[leftmargin=1.5em,nosep]
  \item the test identifier and category;
  \item the result: \emph{pass}, \emph{conditional pass} (with
    documented condition), \emph{fail};
  \item the evidence: log excerpt, command output, screenshot, or
    artefact reference;
  \item the remediation owner and timeline, if applicable.
\end{itemize}

\paragraph{Exit-confidence index per layer.} A summary index, on the
same four-point scale as the sovereignty grid (0 Absent, 1 Partial,
2 Established, 3 Robust), is produced per layer based on the test
results. The index reflects the institution's confidence that the
layer can be replaced within a documented time window: 3 if all
tests in all four categories pass; 2 if interface fidelity and at
least one of data portability or exit rehearsal pass; 1 if only
interface fidelity passes; 0 if interface fidelity is broken.
The index is graded against the institution's own conformance test
set: an institution extends the illustrative reference suite
(Table~\ref{tab:cf-suite}) to cover all four categories for each
layer it operates, so a 3 (Robust) is attainable wherever that
coverage is implemented --- the reference suite is a starting
illustration, not the grading baseline.

\paragraph{The index is a maturity trajectory, not a day-one gate.}
The exit-confidence index measures where the institution stands on its
path to replaceability; it is not a pass/fail compliance verdict and
no institution is expected to reach 3 (Robust) on every layer at the
outset. The starting baseline is diagnosed by the self-assessment of
\trchap{ch:self-assessment}{}, on the same four-point scale, and the
trajectory and pace toward higher confidence are set by the playbook
(\trchap{ch:playbook-reading}{}); the index then records progress
along that trajectory. A smaller or teaching-led institution following
the down-scaled profile (\trchap{ch:sizing}{},
\S\ref{sec:sz-downscale}) may legitimately target a steady-state index
below 3 on layers it consumes as managed or mutualised services: for
such a layer the replaceability question is answered by the provider's
contractual exit terms rather than by a test the institution runs
itself.

%-----------------------------------------------------------------------------
\section{Linkage with wave closure}
\label{sec:cf-linkage}

Every wave of Part~III closes by passing the conformance subset
applicable to its scope. This is the mechanism that prevents the
common failure mode of migration programmes: the change of product
disguised as architectural progress. A wave that delivered a new
identity provider but did not pass the identity-layer
data-portability and exit-rehearsal tests has not closed; the
institution has merely changed supplier.

\paragraph{Wave-to-test mapping.} The mapping is published in each
wave chapter (Chapters~18--27) under the \emph{exit criteria}
heading. The reference mapping is, briefly:
Wave~1 (Identity) closes against the Identity-layer tests; Wave~2
(Observability) against the Observability-layer tests; Wave~3
(Network) against the Network-layer tests; Wave~4 (Policy) against
the Policy-layer tests; Wave~5 (Pilot) against the tests on the
layers exercised by the pilot, across the categories they implement;
Wave~6 (Endpoint) against the Endpoint-layer tests; Wave~7 (Audit and IR)
against an Observability subset reinforced with audit-replay
exercises; Wave~8 (Workload) against the Compute and Storage tests;
Wave~9 (Lifecycle) against the Orchestration-layer tests.
The Crypto layer (Ch.~13) is a cross-cutting foundation rather than a
standalone wave: certificate and key services underpin Identity, so
the layer is stood up as the technical prerequisite of Wave~1
(Identity), and its three tests --- ACME/RFC~8555 issuance,
key-escrow recovery, and HSM migration --- close with Wave~1's exit
criteria. With this, all nine architecture layers are closed across
the roadmap.

\medskip

\noindent The conformance suite is not a one-time effort. Its value
accumulates over time, as the suite reflects the institution's
specific risks, as the test results form a longitudinal record of
the architecture's health, and as the suite itself becomes an
artefact that the institution can share with peers, where it
participates in a peer or research-network community.
